\documentclass[12pt,a4paper]{report}

% Configurações de codificação e idioma
\usepackage[utf8]{inputenc}
\usepackage[T1]{fontenc}
\usepackage[brazilian]{babel}

% Geometria da página
\usepackage[top=3cm,bottom=2cm,left=3cm,right=2cm]{geometry}

% Matemática e teoremas
\usepackage{amsmath,amssymb,amsthm}
\theoremstyle{definition}
\newtheorem{definicao}{Definição}[chapter]
\newtheorem{exemplo}{Exemplo}[chapter]
\theoremstyle{plain}
\newtheorem{teorema}{Teorema}[chapter]
\newtheorem{lema}{Lema}[chapter]
\newtheorem{corolario}{Corolário}[chapter]

% Algoritmos com algorithm2e (rótulos em português)
\usepackage[ruled,linesnumbered,portuguese]{algorithm2e}

% Gráficos, tabelas e diagramas
\usepackage{graphicx}
\usepackage{booktabs}
\usepackage{tikz}
\usepackage{pgfplots}
\pgfplotsset{compat=1.18}
\usetikzlibrary{arrows.meta,positioning,shapes.geometric,trees,matrix}

% Código fonte e formatação
\usepackage{listings}
\usepackage{xcolor}
\lstset{
  basicstyle=\ttfamily\small,
  breaklines=true,
  frame=single,
  numbers=left,
  numberstyle=\tiny,
  keywordstyle=\color{blue},
  commentstyle=\color{gray}
}

% Links e referências cruzadas
\usepackage{hyperref}
\hypersetup{
  colorlinks=true,
  linkcolor=blue,
  citecolor=blue,
  urlcolor=blue
}

% Informações do Documento
\title{
  \vspace{-2cm}
  \textbf{Centro Federal de Educação Tecnológica de Minas Gerais}\\
  \large Departamento de Computação -- Campus Divinópolis\\[1.5cm]
  \Large \textbf{Estratégias Avançadas de Projeto e Análise de Algoritmos}\\
  \huge \textbf{Programação Dinâmica}\\
  \large Fundamentação Teórica, Implementação de Referência e Avaliação Empírica\\[2cm]
}

\author{
  \textbf{Equipe de Desenvolvimento:}\\[0.3cm]
  [Nome do Integrante 1] -- Relatório (Introdução, Fundamentação, Conclusão)\\
  [Nome do Integrante 2] -- Relatório (Modelagem, Pseudocódigos, Análise Formal)\\
  [Nome do Integrante 3] -- Implementação Base dos Algoritmos em Python\\
  [Nome do Integrante 4] -- Estudo de Caso, Benchmarks, Gráficos e C++\\
  [Nome do Integrante 5] -- Cheat Sheet, Slides e Visualizadores Web\\[1.5cm]
  \textbf{Professor Orientador:} Prof. Michel Pires da Silva\\[0.3cm]
  \textbf{Disciplina:} Algoritmos e Estruturas de Dados II
}

\date{Divinópolis -- MG\\ 2026}

\begin{document}

\maketitle

\tableofcontents
\newpage

% Inclusão sequencial das seções
\chapter{Introdução}
\label{cap:introducao}

% TODO(P1): Estruturar a introdução no modelo em funil clássico:
% 1. Contextualização: importância dos problemas de otimização combinatória na computação e indústria real.
% 2. O Problema / Limitação: ineficiência de algoritmos recursivos ingênuos que sofrem com explosão combinatória
%    devido à recomputação desnecessária de subproblemas sobrepostos.
% 3. Objetivo do Trabalho: apresentar, formalizar e validar empiricamente o paradigma de Programação Dinâmica (Cap. 15 do Cormen).
% 4. Entregáveis do Grupo: relatório técnico-científico, prova de conceito modular em Python/C++, suíte de experimentos e simuladores web.
% 5. Roteiro de Leitura: breve resumo do que cada capítulo deste relatório aborda (Fundamentação, Paradigma, Análise, Estudo de Caso e Conclusões).

% [Texto a ser redigido pelo integrante P1 - ver docs/tarefas/P1_relatorio_secoes_1_2_6.md]

\chapter{Fundamentação Teórica}
\label{cap:fundamentacao}

% TODO(P1): Redigir a fundamentação teórica formal cobrindo os seguintes itens do Cap. 15 do Cormen:
% 1. Definição formal de Programação Dinâmica como método de projeto de algoritmos para problemas de otimização.
% 2. Propriedade 1: Subestrutura Ótima.
%    - Apresentar a demonstração formal pela técnica de "cortar e colar" (cut-and-paste) por contradição,
%      aplicada ao problema do Corte de Hastes (Cormen 15.1).
% 3. Propriedade 2: Superposição de Subproblemas (Overlapping Subproblems).
%    - Contrastar com divisão e conquista (onde os subproblemas são independentes).
%    - Grafo de subproblemas direcionado (DAG) e equivalência assintótica do número de vértices/arestas.
% 4. Os Quatro Passos do Projeto segundo Cormen:
%    - Passo 1: Caracterizar a estrutura de uma solução ótima.
%    - Passo 2: Definir recursivamente o valor de uma solução ótima.
%    - Passo 3: Calcular o valor de uma solução ótima (Bottom-Up vs. Top-Down com Memoização).
%    - Passo 4: Construir a solução ótima a partir das informações calculadas (reconstrução/traceback).
% 5. Comparativo Teórico: Programação Dinâmica vs. Algoritmos Gulosos vs. Divisão e Conquista.
% 6. Explosão Combinatória e Números de Catalan:
%    - Apresentar a fórmula P(n) = C(n-1) = (1/n) * \binom{2n-2}{n-1} para a cadeia de matrizes.
%    - Provar que P(n) = \Omega(4^n / n^{3/2}), justificando a necessidade incontornável da PD.

% [Texto a ser redigido pelo integrante P1 - ver docs/tarefas/P1_relatorio_secoes_1_2_6.md]

\chapter{Paradigma Algorítmico e Modelagem}
\label{cap:paradigma}

% TODO(P2): Apresentar os algoritmos completos em pseudocódigo utilizando o pacote algorithm2e
% e os diagramas conceituais em TikZ obrigatórios conforme o enunciado:
%
% 1. Sequência de Fibonacci:
%    - Algoritmo Fib-Naive (recursivo exponencial).
%    - Algoritmo Memoized-Fib (com tabela de cache).
%    - Algoritmo Bottom-Up-Fib (iterativo com vetor).
%
% 2. Corte de Hastes (Cormen 15.1):
%    - Algoritmo Cut-Rod (ingênuo).
%    - Algoritmo Memoized-Cut-Rod e auxiliar Memoized-Cut-Rod-Aux.
%    - Algoritmo Bottom-Up-Cut-Rod.
%    - Algoritmo Extended-Bottom-Up-Cut-Rod (retornando r e s).
%    - Algoritmo Print-Cut-Rod-Solution (reconstrução iterativa).
%
% 3. Multiplicação em Cadeia de Matrizes (Cormen 15.2):
%    - Algoritmo Matrix-Chain-Order (preenchimento por diagonais com laço triplo).
%    - Algoritmo Print-Optimal-Parens (reconstrução recursiva da parentização).
%
% 4. Subsequência Comum Máxima - LCS (Cormen 15.4):
%    - Algoritmo LCS-Length (tabelas c e b).
%    - Algoritmo Print-LCS (traceback da subsequência ótima).
%
% 5. Diagramas Obrigatórios em TikZ (em relatorio/tikz/):
%    - Diagrama (a): Árvore de recursão do Cut-Rod para n = 4 com subproblemas repetidos coloridos.
%    - Diagrama (b): Matriz de programação dinâmica do LCS preenchida com setas direcionais ('diag', 'up', 'left').
%    - Diagrama (c): Grafo direcionado acíclico (DAG) de dependências dos subproblemas.
%    - Diagrama (d): Tabelas m e s completas para o exemplo de 6 matrizes do Cormen (dims = [30, 35, 15, 5, 10, 20, 25]).
%    - Diagrama (e - opcional): Estado da pilha de execução e cache durante a memoização.

% [Texto, pseudocódigos e diagramas TikZ a cargo de P2 - ver docs/tarefas/P2_relatorio_secoes_3_4.md]

\chapter{Análise Assintótica e Complexidade}
\label{cap:analise}

% TODO(P2): Apresentar análise matemática rigorosa de tempo e espaço com notações O, Omega e Theta:
%
% 1. Método Geral de Análise em Programação Dinâmica:
%    - Tempo Total = (Número total de subproblemas distintos) * (Custo de decisão/combinação por subproblema).
%    - Espaço Total = (Tamanho da tabela de armazenamento) + (Profundidade da pilha de recursão no Top-Down).
%
% 2. Prova Formal de Ineficiência do Corte de Hastes Ingênuo:
%    - Recorrência: T(0) = 1 e T(n) = 1 + \sum_{j=0}^{n-1} T(j).
%    - Demonstração por indução matemática de que T(n) = 2^n = \Theta(2^n).
%
% 3. Análise dos Algoritmos Otimizados:
%    - Fibonacci: Naive \Theta(\phi^n) vs. PD \Theta(n) tempo e \Theta(1) espaço.
%    - Corte de Hastes: Bottom-Up \Theta(n^2) tempo e \Theta(n) espaço.
%    - Cadeia de Matrizes: \Theta(n^3) tempo e \Theta(n^2) espaço (laço triplo sobre subcadeias e pontos de partição).
%    - LCS: \Theta(m \cdot n) tempo e \Theta(m \cdot n) espaço para tabela completa, reduzido a \Theta(\min(m, n)) espaço com duas linhas.
%
% 4. Tabela-Resumo Oficial de Complexidades:
%    - Construir tabela com colunas: Problema | Abordagem | Tempo (Pior/Médio) | Espaço | Observações.
%    - Esta tabela será espelhada diretamente no Cheat Sheet oficial.

% [Texto e deduções matemáticas a cargo de P2 - ver docs/tarefas/P2_relatorio_secoes_3_4.md]

\chapter{Estudos de Caso e Resultados Experimentais}
\label{cap:estudo_de_caso}

% TODO(P4): Detalhar a implementação prática, os estudos de caso reais e a discussão crítica dos benchmarks:
%
% 1. Estudos de Caso Práticos:
%    - Ferramenta de Diff Textual (diff_tool.py): aplicação direta de LCS considerando linhas como tokens,
%      comportamento em arquivos reais de código/texto em data/.
%    - Bioinformática e Alinhamento de DNA (dna_alignment.py): distância de edição de Levenshtein e
%      alinhamento global com inserção de gaps ('-') sobre arquivos FASTA.
%
% 2. Metodologia Experimental:
%    - Instrumentação de tempo (time.perf_counter) e memória (tracemalloc).
%    - Protocolo com mediana de 5 repetições, semente aleatória fixa e corte de 10s para algoritmos exponenciais.
%    - Metadados do ambiente de execução (hardware, CPU, sistema operacional, versões do Python e g++).
%
% 3. Análise dos Resultados Gráficos (inserir figuras geradas por plot_results.py em relatorio/figuras/):
%    - Experimento E1: Fibonacci (tempo e chamadas vs. n em escala logarítmica).
%    - Experimento E2: Corte de Hastes (\Theta(2^n) vs. \Theta(n^2)).
%    - Experimento E3: LCS - consumo de memória (tabela completa vs. duas linhas).
%    - Experimento E4: Gargalo de Linguagem - Python vs. C++17 compilado com -O3.
%    - Experimento E5: Limite de Recursão do Python (sys.getrecursionlimit() na memoização).
%    - Experimento E6: Taxa de falha empírica da heurística gulosa no corte de hastes.
%    - Experimento E7: Cadeia de Matrizes - tempo da PD vs. crescimento dos números de Catalan.
%
% 4. Discussão Crítica de Gargalos:
%    - Overhead do interpretador Python frente ao C++.
%    - Impacto da localidade espacial de cache (acesso a matrizes por linhas vs. colunas).
%    - Consumo excessivo de memória na matriz 2D em strings longas.

% [Texto e figuras a cargo de P4 - ver docs/tarefas/P4_estudo_de_caso_e_benchmarks.md]

\chapter{Conclusões e Trabalhos Futuros}
\label{cap:conclusoes}

% TODO(P1): Sintetizar as conclusões finais do estudo, cobrindo:
%
% 1. Síntese dos Resultados:
%    - Confirmação empírica do ganho assintótico da Programação Dinâmica em relação a abordagens ingênuas.
%    - Papel essencial do reconhecimento prévio da subestrutura ótima e sobreposição de subproblemas.
%
% 2. Onde a Programação Dinâmica Falha ou é Limitada:
%    - Ausência de Subestrutura Ótima: problema do caminho simples mais longo em grafos (subproblemas não são independentes).
%    - Explosão do Espaço de Estados: Caixeiro Viajante (TSP) com o algoritmo de Held-Karp atingindo O(n^2 \cdot 2^n).
%    - Natureza Pseudo-Polinomial: Problema da Mochila 0/1 com complexidade O(n \cdot W), onde W é o valor numérico.
%    - Gargalo de Memória: problemas com matrizes multidimensionais onde o espaço Theta(n^2) ou Theta(n^3) esgota a RAM.
%
% 3. Trabalhos Futuros:
%    - Paralelização de preenchimento de tabelas em GPUs/multithreading.
%    - Heurísticas aproximativas para problemas NP-difíceis onde a PD clássica é inviável.

% [Texto a cargo de P1 - ver docs/tarefas/P1_relatorio_secoes_1_2_6.md]


% Apêndices
\appendix
\chapter{Cheat Sheet: Guia Rápido de Programação Dinâmica}
\label{apendice:cheatsheet}

% Inclui o corpo reutilizável do cheat sheet diretamente neste apêndice
% ==============================================================================
% Corpo Reutilizável do Cheat Sheet de Programação Dinâmica
% Compartilhado entre cheatsheet/cheatsheet.tex e relatorio/secoes/apendice_cheatsheet.tex
% Responsável: Integrante P5 (com tabela alimentada por P2)
% ==============================================================================

\section*{Guia Rápido: Programação Dinâmica em Uma Página}

% ------------------------------------------------------------------------------
% Bloco (a): Definição em Uma Frase
% ------------------------------------------------------------------------------
\subsection*{1. Definição Essencial}
% TODO(P5): Inserir a definição canônica concisa em uma única frase.
% Exemplo: "Técnica de projeto para problemas de otimização que resolve cada subproblema
% apenas uma vez, armazenando seu resultado em tabela para reaproveitamento futuro."
\noindent\fbox{%
  \parbox{\dimexpr\linewidth-2\fboxsep-2\fboxrule}{%
    \textbf{Definição:} \textit{[TODO(P5): Inserir a definição em uma frase aqui]}
  }%
}

\vspace{0.5cm}

% ------------------------------------------------------------------------------
% Bloco (b): Diagrama de Decisão "Quando utilizar esta abordagem?"
% ------------------------------------------------------------------------------
\subsection*{2. Diagrama de Decisão: Quando Utilizar?}
% TODO(P5): Completar o fluxograma de decisão em TikZ:
% 1. O problema é de otimização? (Não -> Buscar outra técnica; Sim -> Próximo)
% 2. Possui Subestrutura Ótima? (Não -> PD não se aplica; Sim -> Próximo)
% 3. Os subproblemas se sobrepõem?
%    - Se Não (independentes) -> Divisão e Conquista.
%    - Se Sim -> Uma escolha local gulosa é sempre globalmente ótima?
%      * Se Sim -> Algoritmo Guloso.
%      * Se Não -> PROGRAMAÇÃO DINÂMICA.

\begin{center}
\begin{tikzpicture}[
    node distance=1.5cm,
    decision/.style={diamond, draw, fill=blue!10, text width=5em, text centered, inner sep=0pt, font=\footnotesize},
    block/.style={rectangle, draw, fill=gray!10, text width=6em, text centered, rounded corners, minimum height=2em, font=\footnotesize},
    line/.style={draw, -Latex, font=\scriptsize}
]
  % Nós em esqueleto
  \node [decision] (opt) {Problema de Otimização?};
  \node [decision, right=2cm of opt] (sub) {Subestrutura Ótima?};
  \node [decision, below=1.2cm of sub] (over) {Subproblemas Sobrepostos?};
  \node [block, right=1.5cm of over] (pd) {\textbf{Programação Dinâmica}};
  \node [block, below=1cm of over] (dc) {Divisão e Conquista};
  \node [block, below=1.2cm of opt] (other) {Outras Abordagens};

  % Arestas conceituais
  \path [line] (opt) -- node[above]{Sim} (sub);
  \path [line] (opt) -- node[left]{Não} (other);
  \path [line] (sub) -- node[right]{Sim} (over);
  \path [line] (over) -- node[above]{Sim} (pd);
  \path [line] (over) -- node[right]{Não} (dc);
\end{tikzpicture}
\end{center}

\vspace{0.3cm}

% ------------------------------------------------------------------------------
% Bloco (c): Tabela Oficial de Complexidades
% ------------------------------------------------------------------------------
\subsection*{3. Tabela Comparativa de Complexidades}
% TODO(P5/P2): Preencher a tabela com os valores consolidados da Seção 4.
\begin{table}[h!]
\centering
\small
\begin{tabular}{@{}lllll@{}}
\toprule
\textbf{Problema} & \textbf{Ingênuo (Tempo)} & \textbf{PD (Tempo)} & \textbf{PD (Espaço)} & \textbf{Recorrência Chave} \\ \midrule
Fibonacci & $\Theta(2^n)$ & $\Theta(n)$ & $\Theta(n)$ ou $\Theta(1)$ & $F(n) = F(n-1) + F(n-2)$ \\
Corte de Hastes & $\Theta(2^n)$ & $\Theta(n^2)$ & $\Theta(n)$ & $r_n = \max (p_i + r_{n-i})$ \\
Cadeia de Matrizes & $\Omega(4^n / n^{3/2})$ & $\Theta(n^3)$ & $\Theta(n^2)$ & $m[i,j] = \min (m[i,k]+m[k+1,j]+p_{i-1}p_kp_j)$ \\
LCS & $O(2^{m+n})$ & $\Theta(m \cdot n)$ & $\Theta(m \cdot n)$ ou $\Theta(\min)$ & $c[i,j] = c[i-1,j-1]+1 \text{ ou } \max$ \\
Mochila 0/1 (Opcional) & $O(2^n)$ & $O(n \cdot W)$ & $O(n \cdot W)$ ou $O(W)$ & $V[i,w] = \max(V_{sem}, V_{com} + v_i)$ \\ \bottomrule
\end{tabular}
\end{table}

\vspace{0.3cm}

% ------------------------------------------------------------------------------
% Bloco (d): Esqueleto de Pseudocódigo (Top-Down vs. Bottom-Up)
% ------------------------------------------------------------------------------
\subsection*{4. Padrões de Implementação: Top-Down vs. Bottom-Up}
% TODO(P5): Inserir o template genérico dos dois estilos clássicos de PD.

\noindent
\begin{minipage}[t]{0.48\textwidth}
\textbf{Top-Down (com Memoização):}
\begin{algorithm}[H]
\caption{Padrao-Top-Down(x, memo)}
\SetKwFunction{FPD}{Padrao-Top-Down}
\If{caso base atingido}{
  \Return valor base\;
}
\If{x já está em memo}{
  \Return memo[x]\;
}
res $\leftarrow$ combinar escolhas recursivas\;
memo[x] $\leftarrow$ res\;
\Return res\;
\end{algorithm}
\end{minipage}
\hfill
\begin{minipage}[t]{0.48\textwidth}
\textbf{Bottom-Up (Tabulação Iterativa):}
\begin{algorithm}[H]
\caption{Padrao-Bottom-Up(n)}
Alocar tabela T[0..n]\;
Inicializar casos base em T\;
\For{tamanho $\leftarrow 1$ \KwTo $n$}{
  T[tamanho] $\leftarrow \min/\max$ sobre subproblemas já calculados\;
}
\Return T[n]\;
\end{algorithm}
\end{minipage}



% Bibliografia
\bibliographystyle{plain}
\bibliography{referencias}

\end{document}
